\documentclass{article}
\usepackage{graphicx} % Required for inserting images

\title{Projekt-Praktikum}
\author{Javier Bellido Roldan}
\date{March 2026}

\begin{document}

\maketitle

\section{Introduction}
The rotational dynamics of thin rigid disks (popularly know as Euler's disks) has been a subject of study for many centuries (Euler, 18th Century). In real systems, the energy of a system is gradually dissipated until the object in motion comes to rest. Particularly in the case of an Euler disk under rotation in a plane, the precession frequency of the object gradually approaches infinity as the energy is dissipated. Although the analytical equations of motion are well known, the determination of the dominant dissipative mechanism has been in debate in recent times (and a ArXiV preprint from 15. March 2026, by researchers from Harvard University, challenges the current dominant thesis!). Most studies on the subject rely on computational methods for evaluation, but spectral analysis techniques have been proven sufficiently efficient at accurately measuring its oscillation dynamics. Therefore, we would like to propose an experimental spectral analysis of the rotational dynamics of an Euler disk using inexpensive equipment. Moreover, we anticipate that such an experiment could be proven useful for analysing the dominant mechanism for the energy dissipation of the system.


\end{document}
